\clearpage\SourceLeaf{196}%
\pdfbookmark[1]{Gauss an Schumacher, 1. April 1843}{schumacher18430401}
\fancyhead[L]{\small GAUSS AN SCHUMACHER}
\begin{center}
{\small GAUSS AN SCHUMACHER.\par}
\medskip\hrule\medskip
{\small Astronomische Nachrichten herausgegeben von H.\ C.\ \textsc{Schumacher}.\par
Band 20.\ Nr.\ 474.\ S.\ 299.\par}
\medskip\hrule
\end{center}\medskip
\begin{flushright}Göttingen, 1843 April 1.\end{flushright}
Um aus Elementen für eine gegebene Zeit einen Ort zu bestimmen, brauche ich zur Berechnung der Anomalie gern die \textsc{Burckhardt}'sche Tafel, die aber nur bis $163^0\ 45'$ geht, und daher für den gegenwärtigen Stand des Cometen nach Herrn \textsc{Galles} Elementen unzureichend wird. \textsc{Barkers} Tafel reicht zwar überall aus, wird aber bei grossen Anomalien wegen des beschwerlichen Interpolirens sehr unbequem. In solchen Fällen pflege ich ein besonderes Verfahren anzuwenden, dessen Mittheilung Ihnen vielleicht angenehm sein wird. Ist $M$ die Zahl mit der (oder für grössere Werthe mit deren Logarithmen) man in die \textsc{Barker}'sche Tafel eingehen müsste, also $M=\frac{\text{Zwischenzeit}}{nq^{\frac32}}$ wo $\log n=0{,}0398723$, so setze ich $\log\frac{MM}{16875}=3P$, und suche in meiner kleinen Logarithmentafel, $A$ und $B$ in der dortigen Bedeutung genommen, der Gleichung $3A+2B=3P$ Genüge zu leisten, was immer, wenn $P$ gross ist, sehr schnell bewirkt wird. Ist dann $a$ die zum Logarithmen $A$ gehörige Zahl, so wird, die Anomalie $=v$ gesetzt,
\[
\operatorname{tang}\tfrac12v=\BareRoot{(3a)}\quad\text{oder}\quad
\log\operatorname{tang}\tfrac12v=\tfrac12(A+\log3)
\]

Auch der Logarithme des Radius Vector wird dann. äusserst bequem berechnet, indem man mit $A+\log3$ wieder in die erste Columne eingeht, oder $A+\log3=A^*$ und die dazu gehörige Grösse in der zweiten Columne $=B^*$ folgt, wodurch sogleich der Logarithme des Radius Vectors $=A^*+B^*+\log q$ wird.

\SourceLeaf{197}%
Die indirecte Auflösung jener Gleichung geschieht, wenigstens für die ersten Versuche, etwas bequemer und fast à vue in der Form $C=P+\frac13B$; man kann zuerst $P$ in der dritten Columne aufsuchen, oder $P=C'$ und die dazu gehörige Grösse in der zweiten Columne $=B'$ setzen, dann $P+\frac13B'=C''$ und dazu aus der Tafel die Grösse der zweiten Columne $=B''$, dann (wo nöthig) $P+\frac13B''=C'''$ und dazu gehörig $B'''$ nehmen u.\ s.\ w., welche Rechnung sehr schnell zum Stillstand kommt. Will man sich mit der Genauigkeit, welche fünfzifrige Logarithmen geben, nicht begnügen, so kann man die \textsc{Matthiessen}'sche Tafel (welche ich sonst wegen der unzeitigen Oeconomie, womit sie ganz unnöthigerweise gedruckt ist, nicht gern gebrauche) hier mit Vortheil zu Hülfe nehmen, was ich aber lieber erst dann thue, wenn ich durch die kleinere Tafel die beiden Stellen, zwischen welchen der Definitivwerth von $A$ fällt schon bestimmt habe, und dann wende ich lieber die Gleichung in ihrer ursprünglichen Form $3A+2B=3P$ an.

Soll z.\ B.\ die Anomalie für Februar $48.33333$, oder für die Zeit nach der Sonnennähe $20^{\mathrm T}87663$ bestimmt werden, so ist nach \textsc{Galle}'s Elementen [Astr.\ Nachrichten Band 20.\ N.\ 474 Berlin 1843 März 25]
\begin{center}
{\small\setlength{\tabcolsep}{3pt}
\begin{tabular}{r@{\ }r@{\qquad}r@{\ }r}
$q^{\frac32}\ .\ .\ .\ .\ .\ .$&$7.0809490$&$20.87663\ .\ .\ .\ .$&$1.3196604$\\
$n\BareRoot{16875}\ .\ .$&$2.1534942$&&$9.2344432$\\\cline{1-2}\cline{4-4}
Const.\ Logarithme $=$&$9.2344432$&&$2.0852172$\\[3pt]
&&Also $3P=$&$4.1704344$\\
&&$P=$&$1.3901448$
\end{tabular}}
\end{center}
Mit den kleinen Tafeln findet sich daraus
\[
\begin{array}{ll}
B'=0.01806&\quad C''=1.39616\\
B''=0.01781&\quad C'''=1.39608
\end{array}
\]
womit die Rechnung schon steht, und $A=1.37827$ wird. \textsc{Matthiessen}'s Tafel gibt genauer $A=1.3782739$. Die weitere Rechnung wird dann
\begin{center}
\begin{tabular}{r@{\ }l}
$A=$&$1.3782739$\\
$3\ .\ .\ .\ .\ .$&$0.4771213$\\\cline{2-2}
&$1.8553952$\\[3pt]
&$0.9276976=\log\operatorname{tang}\ 83^0\ 15'\ 49''53$\\
\multicolumn{2}{r}{und die wahre Anomalie $=\qquad166\ 31\ 39{,}06$}
\end{tabular}
\end{center}

\SourceLeaf{198}%
Ferner gehört zu
\begin{center}
\begin{tabular}{r@{\ }r}
$A^*=$&$1.8553952$\\
$B^*=$&$0.0060170$\\
$q\ .\ .\ .\ .$&$8.0539660$\\\cline{1-2}
Logarithm des Radius Vector $=$&$9.9153782$
\end{tabular}
\end{center}

Man sieht übrigens, dass diese Methode nichts weiter ist, als eine indirecte Auflösung der bekannten cubischen Gleichung zwischen der Tangente der halben Anomalie und der Sectorfläche, und zugleich, dass meine, oder für schärfere Rechnung die \textsc{Matthiessen}'sche Logarithmentafel auf ganz ähnliche Weise zu einer sehr bequemen Auffindung aller reellen Wurzeln \textit{jeder algebraischen Gleichung, die nicht mehr als drei effective Glieder hat}, benutzt werden kann, wie ich in Beziehung auf die quadratische Gleichung unlängst bei der letzten Ausgabe der \textsc{Vega}'schen Logarithmentafel schon gezeigt habe.
\par\medskip\centerline{\rule{.25\textwidth}{.3pt}}
